\chapter{Conclusiones y Trabajo Futuro}

\section{Síntesis de Hallazgos}

Esta tesis se propuso comparar, bajo condiciones de rigurosa imparcialidad metodológica, el desempeño de los métodos de interpolación de canales EEG basados en Procesamiento de Señales en Grafos con regularización temporal frente a los \textit{baselines} geométricos clásicos. Tras un proceso iterativo de implementación, experimentación exploratoria, optimización bayesiana y validación confirmatoria, los hallazgos principales pueden sintetizarse en cinco conclusiones.

\textbf{1. La regularización temporal marca una diferencia cualitativa.} TRSS no solo reduce el error de amplitud frente a los métodos puramente espaciales —incluso cuando estos son exhaustivamente optimizados—, sino que fundamentalmente altera el \textit{modo de falla} de la interpolación. Donde los métodos geométricos colapsan (inversión de fase, divergencia de amplitud), TRSS mantiene trayectorias fisiológicamente plausibles. Esta diferencia no es de grado sino de naturaleza: la restricción temporal cambia la topología del espacio de soluciones, eliminando los mínimos locales espurios que atrapan a los métodos muestra-a-muestra.

\textbf{2. La contribución del término temporal es asimétrica.} El estudio de ablación reveló que $\beta$ no mejora significativamente el error punto-a-punto (MAE, RMSE), pero produce mejoras grandes y estadísticamente significativas en fidelidad espectral (LSD: $+8\%$, Cliff's $\delta = -0.54$ a $-0.76$, $p < 0.01$). Esta disociación tiene una interpretación física clara: el error de amplitud está gobernado por la estructura espacial del grafo; la preservación frecuencial depende de la continuidad temporal. Un método de interpolación completo requiere ambos.

\textbf{3. MNE interpolate\_bads es un baseline excepcionalmente fuerte que no debe subestimarse.} La implementación industrial del \textit{spherical spline} en MNE-Python incorpora décadas de refinamientos ingenieriles —normalización de posiciones, regularización diagonal, aumento del sistema, pseudoinversa con tolerancia adaptativa— que lo convierten en un contrincante formidable. En escenarios favorables (montajes densos, $\leq$20\% de pérdida aleatoria), MNE iguala o supera a TRSS con un costo computacional 50--100$\times$ menor. Cualquier trabajo futuro que proponga un nuevo método de interpolación debe incluir a MNE en su panel de comparación.

\textbf{4. La recomendación de método es condicional al contexto.} No existe un método universalmente óptimo. TRSS se justifica cuando la prioridad es la preservación de dinámica temporal fina (ERPs, conectividad funcional, análisis espectral) bajo condiciones de pérdida severa ($\geq$30\%) o contigua. MNE interpolate\_bads es la opción correcta para escenarios de rutina con restricciones de tiempo, montajes densos y patrones de pérdida aleatorios de baja severidad.

\textbf{5. El \textit{fair benchmarking} es posible y necesario.} La optimización bayesiana de \textit{todos} los métodos —incluyendo los baselines clásicos— con separación cronológica estricta entre datos de optimización y prueba establece un estándar metodológico que esta tesis espera ver adoptado en la literatura de interpolación de señales biomédicas. Los resultados obtenidos bajo este protocolo son más conservadores —las ventajas reportadas son menores— pero también más confiables.

\section{Cumplimiento de Objetivos}

El objetivo general —comparar el rendimiento de métodos GSP con regularización temporal frente a baselines geométricos mediante un marco experimental riguroso— ha sido plenamente alcanzado. Los siete objetivos específicos definidos en la Sección~1.5 se cumplieron de la siguiente manera:

\begin{enumerate}
    \item \textbf{Unificación metodológica:} Se implementó un \textit{pipeline} modular en Python con 18 métodos de interpolación, 8 constructores de grafo y 6 métricas de evaluación, todos con interfaces homogéneas (Capítulo~3).
    \item \textbf{Escenarios clínicamente representativos:} Se definieron 14 escenarios de pérdida combinando modos espaciales (random, nearby) con niveles de severidad del 10\% al 40\% (Sección~3.4).
    \item \textbf{Optimización imparcial:} Se aplicó Optuna con TPE y NSGA-II a los 5 métodos seleccionados, con espacios de búsqueda documentados y protocolo de validación temporal (Sección~3.6).
    \item \textbf{Evaluación multi-métrica:} Se cuantificó el desempeño en 6 métricas sobre 3 datasets, con 736 iteraciones y 4336 puntos de dato (Capítulo~4).
    \item \textbf{Relevancia fisiológica:} Se demostró la preservación de ERPs (N100, P300) y PSD en bandas clínicas mediante figuras comparativas (Sección~4.3 y~4.4).
    \item \textbf{Validación confirmatoria contra MNE:} Se realizó una comparación directa TRSS vs. MNE interpolate\_bads con máscara estática y optimización independiente (Sección~\ref{sec:resultados_mne}).
    \item \textbf{Trazabilidad y reproducibilidad:} El código fuente, las configuraciones experimentales y los artefactos de resultados están disponibles en un repositorio Git público.
\end{enumerate}

\section{Trabajo Futuro}

Los resultados de esta tesis habilitan varias direcciones de investigación que se consideran prioritarias.

\textbf{Validación clínica.} Extender la evaluación a registros de pacientes con patologías neurológicas (epilepsia, enfermedad de Alzheimer, trastornos de la conciencia) donde la señal EEG presenta características anormales (puntas epilépticas, enlentecimiento difuso) que pueden desafiar los \textit{priors} de suavidad asumidos por todos los métodos evaluados.

\textbf{Aprendizaje de hiperparámetros.} El acoplamiento entre $\alpha$, $\beta$ y la frecuencia de muestreo $f_s$ fue resuelto por Optuna en este trabajo, pero una caracterización sistemática de esta relación —potencialmente derivable analíticamente— permitiría transferir hiperparámetros entre datasets sin re-optimización.

\textbf{Aceleración computacional.} Explorar implementaciones en GPU del descenso de gradiente de TRSS, \textit{early stopping} basado en la norma del gradiente, y precondicionadores que reduzcan el número de iteraciones requeridas, para cerrar la brecha de velocidad con MNE.

\textbf{Generalización entre sujetos.} Evaluar la transferibilidad de hiperparámetros optimizados en un sujeto a otros sujetos dentro del mismo dataset, y entre datasets con distintas densidades de electrodos.

\textbf{Comparación con métodos de aprendizaje profundo.} Incluir autoencoders convolucionales y GNNs en el \textit{benchmark}, con volúmenes de entrenamiento controlados, para cuantificar el \textit{trade-off} entre desempeño, requerimiento de datos y costo computacional entre enfoques algorítmicos y basados en datos.

\textbf{Extensión a otros dominios.} La formulación de TRSS es independiente del dominio de aplicación. Su evaluación en otros problemas de señales sobre grafos con pérdida de nodos —redes de sensores industriales, datos climáticos, señales de tráfico— podría revelar patrones de desempeño generalizables.

\section{Comentario Final}

Esta tesis no propone un nuevo algoritmo, sino una nueva forma de evaluar algoritmos existentes. En un campo donde la literatura está plagada de comparaciones incompletas y \textit{benchmarks} sesgados, el valor de establecer un estándar de evaluación riguroso, reproducible y justo puede ser, a largo plazo, más significativo que el de proponer una fórmula de regularización incremental. Los métodos de interpolación de EEG son una herramienta; esta tesis aspira a ser un manual de uso que indique, con evidencia, cuándo y por qué utilizar cada una.